%% BLOCK 4 -- Managing ML Projects.
%% Fragment of ai_foundations.tex; not a standalone document.
%% Source: developers.google.com/machine-learning/managing-ml-projects, all
%% eleven pages (Overview, Development phases, Assembling a team, Working with
%% stakeholders, Feasibility, Planning, Measuring success, Experiments, ML
%% pipelines, Productionization, AI and ML ethics, ML resources).

\actpage{4}{Managing ML Projects}{From an idea to a production-ready implementation: the phases, the team, the plan, and the pipelines that keep a model fresh.}

\begin{frame}{Why traditional ML, in a generative age}
  \begin{withmargin}
    \begin{tdmain}
      This course focuses on traditional ML models. Although generative AI is in
      the spotlight, traditional ML plays a vital role at Google, underpinning
      many services and projects: predicting travel times in Maps, estimating
      the price of airline tickets in Flights, predicting compute quota for
      Cloud customers, recommending relevant videos in YouTube.
      \vskip0.8em
      In general, the principles for managing traditional ML projects are
      \gterm{identical} for managing generative AI projects. Where there's a
      significant difference, the course calls it out.
    \end{tdmain}
    \begin{tdmargin}
      {\color{tdfg}\textbf{Objectives.}}
      \par\medskip
      Define the phases and elements of an ML project. Describe how to plan and
      manage one. Determine business and model success metrics. Recognize the
      iterative process of running ML experiments. Design a solution for
      productionizing ML pipelines. Implement responsible ML and AI practices at
      each development phase.
    \end{tdmargin}
  \end{withmargin}
  \slidesources{GoogleManagingML2025}
\end{frame}

% ======================================================================
\sectionsubtitle{Four phases, each with a goal, tasks, and an outcome.}
\section[Phases]{Development Phases}

\begin{frame}{The four phases}
  \begin{withmargin}
    \begin{tdmain}
      \footnotesize
      \begin{tabular}{@{}p{0.24\linewidth}p{0.32\linewidth}p{0.34\linewidth}@{}}
        \toprule
        \textbf{Phase} & \textbf{Goal} & \textbf{Outcome}\\
        \midrule
        Ideation and planning & Determine if ML is the best solution to your
          problem & A design doc outlining how to solve a problem with an ML
          solution\\[0.4em]
        Experimentation & Build a model that solves the business problem
          & A model with good enough quality to put into production\\[0.4em]
        Pipeline building and productionization & Build and implement the
          infrastructure for scaling, monitoring, and maintaining models in
          production & Validated ML pipelines\\
        \bottomrule
      \end{tabular}
    \end{tdmain}
    \begin{tdmargin}
      A clear understanding of the phases helps establish engineering
      responsibilities, manage stakeholder expectations, and efficiently
      allocate resources.
      \par\medskip
      Successfully moving through the phases, \gterm{often iteratively}, is
      foundational for designing, assembling, and building ML models that solve
      business problems over the long term.
    \end{tdmargin}
  \end{withmargin}
\end{frame}

\begin{frame}{What each phase is underestimated for}
  \begin{withmargin}
    \begin{tdmain}
      \textbf{Ideation and planning.} It can take a long time to understand the
      data and metrics required for a production-ready system.
      \vskip0.6em
      \textbf{Experimentation.} New-to-ML practitioners often underestimate the
      challenges of designing and implementing the appropriate experimentation
      \gterm{tooling and processes}. It's not uncommon to try hundreds of
      experiments before finding the right combination of features,
      hyperparameters, and model architecture.
      \vskip0.6em
      \textbf{Productionization.} It's easy to underestimate the complexity of
      productionizing data pipelines and evaluating models, especially as
      features evolve. Not only do you have to implement all the monitoring
      infrastructure required for a non-ML project, but also all the ML-specific
      monitoring.
    \end{tdmain}
    \begin{tdmargin}
      Multiple challenges exist at each phase. Not realizing, and planning for,
      them may lead to missed deadlines, frustrated engineers, and failed
      projects.
      \par\medskip
      Before spending time drafting a design doc or writing code, you should
      first verify that ML is the right solution to your problem.
    \end{tdmargin}
  \end{withmargin}
  \keyterms{feature; hyperparameter; model; pipeline; prediction}
\end{frame}

% ======================================================================
\sectionsubtitle{Six roles, and the process docs that align them.}
\section[Team]{Assembling a Team}

\begin{frame}{The six common roles}
  \begin{withmargin}
    \begin{tdmain}
      \footnotesize
      \begin{tabular}{@{}p{0.26\linewidth}p{0.30\linewidth}p{0.34\linewidth}@{}}
        \toprule
        \textbf{Role} & \textbf{Knowledge and skills} & \textbf{Main deliverable}\\
        \midrule
        ML product manager & Deep understanding of ML strengths, weaknesses, and
          the development process; aligns business problems to ML solutions
          & Product requirements document\\[0.3em]
        Engineering manager & Sets, communicates, and achieves team priorities
          & Design docs, project plans, performance evaluations\\[0.3em]
        Data scientist & Quantitative and statistical analysis; identifies and
          tests features, prototypes models & Reports and data visualizations\\[0.3em]
        ML engineer & Designs, builds, productionizes, and manages ML models
          & Deployed model meeting business goals\\[0.3em]
        Data engineer & Builds data pipelines for storing, aggregating, and
          processing data & Fully productionized data pipelines with monitoring\\[0.3em]
        DevOps engineer & Develops, deploys, scales, and monitors serving
          infrastructure & Automated serving, monitoring, testing, alerting\\
        \bottomrule
      \end{tabular}
    \end{tdmain}
    \begin{tdmargin}
      Successful ML projects have teams with each role well represented. In
      smaller teams, individuals will need to handle the responsibilities for
      \gterm{multiple roles}.
      \par\medskip
      Data engineers are responsible for the data across the entire ML
      development process.
    \end{tdmargin}
  \end{withmargin}
\end{frame}

\begin{frame}{Process documentation}
  \begin{withmargin}
    \begin{tdmain}
      Because the roles, tools, and frameworks vary widely in ML development,
      it's critical to establish common practices through excellent process
      documentation. Good process docs answer questions like:
      \vskip0.5em
      \footnotesize
      \begin{itemize}
        \item How is the data generated for the model? How do we examine,
          validate, and visualize it?
        \item How do we modify an input feature or label in the training data?
        \item How do we customize the data generation, training, and evaluation
          pipeline?
        \item How do we obtain testing examples? What metrics will we use to
          judge model quality?
        \item How do we launch our models in production? How will we know if
          something is wrong with our model?
        \item What upstream systems do our models depend on?
      \end{itemize}
    \end{tdmain}
    \begin{tdmargin}
      One engineer might think that just getting the right data is sufficient to
      begin training, while a more responsible engineer will validate that the
      dataset is anonymized correctly and document its metadata and provenance.
      \par\medskip
      Making sure engineers share common definitions for processes and design
      patterns reduces confusion and \gterm{increases the team's velocity}.
    \end{tdmargin}
  \end{withmargin}
\end{frame}

\begin{frame}{Everyone is right from their own perspective}
  \begin{withmargin}
    \begin{tdmain}
      What constitutes ML best practices can differ between companies, teams,
      and individuals. Some team members might consider experimental notebooks
      the main deliverable, while others will want to work in R. Some have a
      passion for software engineering, someone else thinks monitoring is the
      most important thing, yet someone else is aware of good feature
      productionization practices but wants to use Scala.
      \vskip0.8em
      Everyone is right from their own perspective and if steered correctly, the
      mix will be a \gterm{powerhouse}. If not, it can be a mess.
    \end{tdmain}
    \begin{tdmargin}
      Establishing the tools, processes, and infrastructure the team will use
      \textbf{before writing a line of code} can be the difference between the
      project failing after two years or successfully launching a quarter ahead
      of schedule.
    \end{tdmargin}
  \end{withmargin}
\end{frame}

\begin{frame}{Performance evaluations under uncertainty}
  \vfill
  \begin{takeaway}
    A team member's performance should not be\\
    directly connected to the \alertbf{success of the project}.
  \end{takeaway}
  \vfill
  {\small It's not uncommon for team members to spend weeks investigating
  solutions that are ultimately unsuccessful. Even in these cases, their
  high-quality code, thorough documentation, and effective collaboration should
  contribute positively toward their evaluation. Due to the ambiguity and
  uncertainty inherent in ML, people managers need to set clear expectations and
  define deliverables early.}
\end{frame}

\begin{frame}{Stakeholders and deliverables}
  \begin{withmargin}
    \begin{tdmain}
      As early as possible, define your project's stakeholders, the expected
      deliverables, and the preferred communication methods.
      \vskip0.6em
      \textbf{Design doc.} Before you write a line of code, create a doc that
      explains the problem, the proposed solution, the potential approaches, and
      possible risks. It functions as a way to receive feedback and address
      concerns.
      \vskip0.5em
      \textbf{Experimental results.} The record of your experiments with their
      hyperparameters and metrics, the training stack, and saved versions of
      your model at certain checkpoints.
      \vskip0.5em
      \textbf{Production-ready implementation.} A full pipeline for training and
      serving. Document modeling decisions, deployment and monitoring specifics,
      and data peculiarities for future engineers.
    \end{tdmain}
    \begin{tdmargin}
      Some stakeholders might assume that ML is similar to traditional software
      engineering with \gterm{deterministic outcomes}. They might not understand
      why the project's progress is stalled or why a project's milestones are
      non-linear.
      \par\medskip
      To manage expectations, be clear about the complexities, timeframes, and
      deliverables at each stage.
    \end{tdmargin}
  \end{withmargin}
\end{frame}

% ======================================================================
\sectionsubtitle{Data, difficulty, quality, technical requirements, cost.}
\section[Feasibility]{Feasibility}

\begin{frame}{Data availability}
  \begin{withmargin}
    \begin{tdmain}
      While ML might be a theoretically good solution, you still need to
      estimate its real-world feasibility. A solution might technically work but
      be \alertbf{impractical or impossible to implement}.
      \vskip0.7em
      \textbf{Quantity.} Can you get enough high-quality data? Are labeled
      examples scarce, difficult to get, or too expensive? Getting labeled
      medical images or translations of rare languages is notoriously hard.
      \vskip0.5em
      \textbf{Feature availability at serving time.} Teams have spent
      considerable amounts of time training models only to realize that some
      features didn't become available until days after the model required them.
      \vskip0.5em
      \textbf{Regulations.} What are the legal requirements for acquiring and
      using the data?
    \end{tdmain}
    \begin{tdmargin}
      A model predicts whether a customer will click a URL, using
      \texttt{user\_age} as a training feature. When the model serves a
      prediction, \texttt{user\_age} isn't available, perhaps because the user
      hasn't created an account yet.
      \par\medskip
      Classification models require numerous examples for \gterm{every} label.
      If the training dataset contains limited examples for some labels, the
      model can't make good predictions.
    \end{tdmargin}
  \end{withmargin}
\end{frame}

\begin{frame}{Data for generative AI}
  \begin{withmargin}
    \begin{tdmain}
      Pre-trained generative AI models often require curated datasets to excel
      at domain-specific tasks.
      \vskip0.5em
      \footnotesize
      \begin{tabular}{@{}p{0.46\linewidth}p{0.46\linewidth}@{}}
        \toprule
        \textbf{Technique} & \textbf{Number of required examples}\\
        \midrule
        Zero-shot prompting & 0\\
        Few-shot prompting & tens to hundreds\\
        Parameter efficient tuning & hundreds to tens of thousands\\
        Fine-tuning & thousands to tens of thousands, or more\\
        \bottomrule
      \end{tabular}
      \vskip0.4em
      Parameter efficient tuning covers LoRA and prompt tuning. \gterm{The
      quality of the fine-tuning data is more important than the quantity.}
    \end{tdmain}
    \begin{tdmargin}
      Once pre-trained, generative AI models have a \textbf{fixed knowledge
      base}. If content in the model's domain changes often, you'll need a
      strategy to keep it up to date: fine-tuning, retrieval-augmented
      generation, or periodic pre-training.
    \end{tdmargin}
  \end{withmargin}
\end{frame}

\begin{frame}{Problem difficulty}
  \begin{withmargin}
    \begin{tdmain}
      What initially appears to be a plausible approach might actually turn out
      to be an open research question. What seems practical and doable might
      turn out to be unrealistic.
      \vskip0.6em
      \textbf{Has a similar problem already been solved?} If so, it's likely
      you'll be able to use parts of their model to build yours. If you're
      attempting a new problem, expect to invest substantial amounts of time and
      resources: you're basically doing a research project that doesn't have a
      guaranteed solution.
      \vskip0.5em
      \textbf{Is the nature of the problem difficult?} Humans classify the type
      of animal in an image at about \gterm{95\%} accuracy, and handwritten
      digits at about \gterm{99\%} accuracy. So a model to classify animals is
      more difficult than one to classify handwritten digits.
    \end{tdmain}
    \begin{tdmargin}
      \textbf{Are there potentially bad actors?} If people will be actively
      trying to exploit your model, you'll be in a constant race to update it
      before it can be misused. Spam filters can't catch new types of spam when
      someone exploits the model to create emails that appear legitimate.
    \end{tdmargin}
  \end{withmargin}
\end{frame}

\begin{frame}{Generative AI vulnerabilities}
  \begin{withmargin}
    \begin{tdmain}
      \textbf{Input source.} Where will the input come from? Can adversarial
      prompts leak training data, preamble material, database content, or tool
      information?
      \vskip0.7em
      \textbf{Output use.} How will outputs be used? Will the model output raw
      content or will there be intermediary steps that test and verify it's
      appropriate? Providing raw output to plugins can cause a number of
      security issues.
      \vskip0.7em
      \textbf{Fine-tuning.} Fine-tuning with a \gterm{corrupted dataset} can
      negatively affect the model's weights, causing it to output incorrect,
      toxic, or biased content.
    \end{tdmain}
    \begin{tdmargin}
      Fine-tuning requires a dataset that's been verified to contain
      high-quality examples.
    \end{tdmargin}
  \end{withmargin}
\end{frame}

\begin{frame}{Prediction quality and its cost}
  \begin{withmargin}
    \begin{tdmain}
      The required prediction quality depends on the type of prediction.
      Recommending the wrong video might create a bad user experience. Wrongly
      flagging a video as violating a platform's policies might generate support
      costs, or worse, legal fees.
      \vskip0.8em
      Generally, the higher the required prediction quality, the harder the
      problem. Projects often reach diminishing returns as you try to improve
      quality. Increasing a model's precision from 99.9\% to 99.99\% could mean
      a \alertbf{10-times increase} in the project's cost, if not more.
    \end{tdmain}
    \begin{tdmargin}
      {\color{tdfg}\textbf{Generative AI output.}}
      \par\medskip
      Although generative models produce fluent and coherent content, it's not
      guaranteed to be factual. False statements from generative AI models are
      called \gterm{confabulations}. Many use cases still require human
      verification before production use, for example LLM-generated code.
      \par\medskip
      What are the legal, financial, or ethical consequences of biased,
      plagiarized, or toxic content?
    \end{tdmargin}
  \end{withmargin}
\end{frame}

\begin{frame}{Technical requirements}
  \begin{withmargin}
    \begin{tdmain}
      \begin{itemize}
        \item \textbf{Latency.} How fast do predictions need to be served?
        \item \textbf{Queries per second.} What are the QPS requirements?
        \item \textbf{RAM usage.} For training and for serving.
        \item \textbf{Platform.} Online, in a web browser, on a phone or tablet,
          or offline with predictions saved in a table?
        \item \textbf{Interpretability.} Will your product need to answer
          questions like "why was a specific piece of content marked as spam"?
        \item \textbf{Retraining frequency.} When the underlying data changes
          rapidly, frequent or continuous retraining may be necessary, at
          significant cost.
      \end{itemize}
      \vskip0.4em
      \footnotesize
      In most cases, you'll have to \gterm{compromise on a model's quality} to
      adhere to its technical specifications.
    \end{tdmain}
    \begin{tdmargin}
      Explaining ML predictions is hard and an active area of research. Tools
      like Model Cards and the What-if Tool can help you begin to understand a
      model's behavior and design for interpretability.
      \par\medskip
      For generative AI, output size affects latency more than input size.
      Models can parallelize their inputs, but they can only generate outputs
      sequentially.
    \end{tdmargin}
  \end{withmargin}
  \slidesources{Mitchell2019}
\end{frame}

\begin{frame}{Cost}
  \begin{withmargin}
    \begin{tdmain}
      Most ML projects won't be approved if the ML solution is more expensive to
      implement and maintain than the money it generates or saves.
      \vskip0.7em
      \textbf{Human costs.} How many people will it take to go from proof of
      concept to production? ML projects require more people to deploy and
      maintain a production-ready system than to create a prototype.
      \vskip0.5em
      \textbf{Machine costs.} Training, deploying, and maintaining models
      requires lots of compute and memory. You might need TPU quota, data
      pipeline infrastructure, paid data labeling, or data licensing fees.
      \vskip0.5em
      \textbf{Inference cost.} Will the model need to make hundreds or thousands
      of inferences that cost \gterm{more than the revenue generated}?
    \end{tdmain}
    \begin{tdmargin}
      Encountering issues related to any of these topics can make implementing
      an ML solution a challenge, but \textbf{tight deadlines amplify} the
      challenges.
      \par\medskip
      Plan and budget enough time based on the perceived difficulty of the
      problem, and then reserve even more overhead time than you might for a
      non-ML project.
    \end{tdmargin}
  \end{withmargin}
  \keyterms{classification model; confabulation; example; few-shot prompt;
    fine-tuning; inference; interpretability; label; labeled example;
    parameter-efficient tuning; precision; pre-trained model; prompt
    engineering; retrieval augmented generation; training; training set;
    zero-shot prompt}
\end{frame}

% ======================================================================
\sectionsubtitle{Non-linear work, planned with an experimental mindset.}
\section[Planning]{Planning}

\begin{frame}{Project uncertainty}
  \begin{withmargin}
    \begin{tdmain}
      Early-stage planning is difficult because the best approach typically
      isn't apparent when beginning a project. This inherent uncertainty makes
      estimating timelines hard.
      \vskip0.8em
      In a Kaggle competition, the first few weeks drew 350 teams, who raised
      benchmark prediction quality from 35\% to 65\%. Over the next two weeks
      the number of teams grew from 350 to 1400. The best model achieved
      \alertbf{68\%}. Over a thousand teams, each experimenting with a variety
      of data transformations, architectures, and hyperparameters, only achieved
      producing a model with 68\% prediction quality.
    \end{tdmain}
    \begin{tdmargin}
      An example from industry illustrates the non-linearity, where the output
      isn't proportional to the inputs. Two teams took several months to train a
      model to 90\% prediction quality. However, it took several teams
      \gterm{more than five years} to get the model ready for production with
      99.9\% prediction quality.
      \par\medskip
      Production-ready ML is an exploratory process, requiring both a scientific
      and an engineering mindset.
    \end{tdmargin}
  \end{withmargin}
\end{frame}

\begin{frame}{Software engineering plans versus ML plans}
  \begin{withmargin}
    \begin{tdmain}
      In a typical software engineering plan, you define the requirements,
      outline the components, estimate the effort, and schedule the work.
      There's a \gterm{clearly defined path} to a solution. Engineers often know
      with a high degree of certainty the tasks they need to complete, and can
      estimate the work based on similar projects.
      \vskip0.8em
      For most ML projects, you'll find the best solution by experimenting with
      multiple approaches in a trial-and-error process. The optimal
      architecture might be a simple linear model, or a neural net, or possibly
      a decision tree. Only by trying each approach can you discover the best
      solution.
    \end{tdmain}
    \begin{tdmargin}
      By definition, experiments aren't guaranteed to be successful. Because of
      this, it's best to plan using an experimental framework.
      \par\medskip
      Only by attempting to solve the problem can you get a better sense for the
      amount of time and resources a solution might require.
    \end{tdmargin}
  \end{withmargin}
\end{frame}

\begin{frame}{Three planning strategies}
  \begin{withmargin}
    \begin{tdmain}
      \textbf{Time box the work.} Allocate two weeks to determine if you can get
      access to the right kind of data. If you can, designate two more weeks to
      see if a simple model indicates feasibility. If a simple model fails,
      designate two more to try a neural net.
      \vskip0.6em
      \textbf{Scope down the requirements.} If an ML solution appears promising
      but isn't a critical feature, plan a very simple solution this quarter and
      improve it iteratively in subsequent quarters.
      \vskip0.6em
      \textbf{Intern or new hire project.} Directing an intern or a new hire to
      attempt an ML solution can be a good way to begin exploring a new space
      with unknown outcomes.
    \end{tdmain}
    \begin{tdmargin}
      At the end of each timeframe, you'll have more information to determine if
      continuing to apply resources to the problem is worthwhile.
      \par\medskip
      With any strategy, it's wise to \gterm{fail fast}. Attempt approaches with
      the lowest costs, but potentially the highest payoff, first. If the
      approach works, you've found a good solution. If it doesn't, you haven't
      wasted lots of time and resources.
    \end{tdmargin}
  \end{withmargin}
\end{frame}

\begin{frame}{Planning as a dynamic document}
  \begin{withmargin}
    \begin{tdmain}
      Learning to plan multiple ML approaches probabilistically takes time and
      experience. Your project plan might require frequent updates. Think of it
      as a dynamic document in \gterm{constant evolution}.
      \vskip0.7em
      \begin{itemize}
        \item Estimate the cost and chance of success for each approach.
        \item Attempt a portfolio of approaches.
        \item Identify lessons learned and try to improve the system one thing
          at a time.
        \item Plan for failures.
      \end{itemize}
    \end{tdmain}
    \begin{tdmargin}
      When an approach leads to a dead end, rather than being discouraged, team
      members understand that's part of the process. By discussing the inherent
      uncertainty with stakeholders, you're able to set more realistic
      expectations.
      \par\medskip
      Occasionally an early approach leads to a breakthrough. Someone might
      discover a bug in the data generation pipeline or the
      training-validation split.
    \end{tdmargin}
  \end{withmargin}
\end{frame}

% ======================================================================
\sectionsubtitle{Business metrics first, model metrics second.}
\section[Success]{Measuring Success}

\begin{frame}{Business metrics}
  \begin{withmargin}
    \begin{tdmain}
      Business metrics are the most important. They're the reason you're using
      ML: \gterm{you want to improve the business}. Start with quantifiable
      product or business metrics, as granular and focused as possible.
      \vskip0.6em
      \begin{itemize}
        \item Reduce a datacenter's monthly electric costs by 30 percent.
        \item Increase revenue from product recommendations by 12 percent.
        \item Increase click-through rate by 9 percent.
        \item Increase customer sentiment from opt-in surveys by 20 percent.
        \item Increase time on page by 4 percent.
      \end{itemize}
    \end{tdmain}
    \begin{tdmargin}
      If you're not tracking the metric you want to improve, start by
      implementing the infrastructure to do so. Setting a goal to increase
      click-through rate 15\% isn't logical if you're not currently measuring
      click-through rates.
      \par\medskip
      The metric you care about might change. You might optimize revenue from
      recommendations, see revenue increase but subscriptions decrease, and
      reevaluate.
    \end{tdmargin}
  \end{withmargin}
\end{frame}

\begin{frame}{Model metrics}
  \begin{withmargin}
    \begin{tdmain}
      \textbf{Determine a single metric to optimize.} Classification models can
      be evaluated against a variety of metrics. Choosing the best model is
      challenging when different metrics favor different models, so agree on
      \gterm{one} metric to evaluate models against.
      \vskip0.7em
      \textbf{Determine acceptability goals to meet.} These are different from
      evaluation metrics. They refer to goals a model needs to meet to be
      considered acceptable for an intended use case, for example "incorrect
      output is less than 0.1\%" or "recall for the top five categories is
      greater than 97\%".
    \end{tdmain}
    \begin{tdmargin}
      A binary classification model detecting fraudulent transactions might
      optimize \textbf{recall}, correctly identifying fraud most of the time,
      while holding \textbf{precision} at or above a particular value as its
      acceptability goal.
      \par\medskip
      Evaluating generative AI output presents unique challenges. For open-ended
      or creative output, it's more difficult than evaluating traditional ML
      outputs.
    \end{tdmargin}
  \end{withmargin}
\end{frame}

\begin{frame}{Model metrics do not guarantee business metrics}
  \begin{withmargin}
    \begin{tdmain}
      A team develops a churn model with 95\% prediction quality and tests it on
      a small sample of users. Revenue doesn't increase. Customer churn actually
      \alertbf{increases}. Possible explanations:
      \vskip0.5em
      \begin{itemize}
        \item Predictions don't occur early enough to be actionable. The model
          can only predict churn within a seven-day timeframe, which isn't soon
          enough to offer incentives.
        \item Incomplete features. Other factors contributing to churn weren't
          in the training dataset.
        \item Threshold isn't high enough. The model might need 97\% or higher
          to be useful.
      \end{itemize}
    \end{tdmain}
    \begin{tdmargin}
      Typically teams deploy the model to \gterm{1\% of users} and then monitor
      the business metric.
      \par\medskip
      Perform early user testing to prove, and understand, the connection
      between the model's metrics and the business metrics.
      \par\medskip
      Many skilled engineers can create models with impressive metrics. Training
      a good-enough model isn't typically the issue. An ML project can be
      destined for failure from a misalignment between business metrics and
      model metrics.
    \end{tdmargin}
  \end{withmargin}
  \keyterms{AUC; AUC-PR; binary classification; F1 score; metric; precision;
    recall; root mean squared error}
\end{frame}

% ======================================================================
\sectionsubtitle{Testable, reproducible hypotheses.}
\section[Experiments]{Experiments}

\begin{frame}{Running experiments}
  \begin{withmargin}
    \begin{tdmain}
      \textbf{Determine baseline performance.} The baseline acts as a measuring
      stick to compare experiments against. The current non-ML solution can
      provide the first baseline. If none exists, create a model with a simple
      architecture and a few features.
      \vskip0.6em
      \textbf{Make single, small changes.} One change at a time, to the
      hyperparameters, architecture, or features. If the change improves the
      model, that model's metrics become the \gterm{new baseline}.
      \vskip0.5em
      \textbf{Record the progress of the experiments.} Experiments with poor or
      neutral prediction quality are still useful to track: they signal which
      approaches won't work.
    \end{tdmain}
    \begin{tdmargin}
      Examples of a single, small change: include feature X; use 0.5 dropout on
      the first hidden layer; take the log transform of feature Y; change the
      learning rate to 0.001.
      \par\medskip
      Because each full training on a real-world dataset can take hours or days,
      consider running multiple independent experiments concurrently to explore
      the space quickly.
    \end{tdmargin}
  \end{withmargin}
\end{frame}

\begin{frame}{Noise in experimental results}
  \begin{withmargin}
    \begin{tdmain}
      You might encounter noise in experimental results that isn't from changes
      to the model or the data, making it difficult to determine if a change
      actually improved the model.
      \vskip0.6em
      \begin{itemize}
        \item \textbf{Data shuffling.} The order in which the data is presented
          can affect performance.
        \item \textbf{Variable initialization.} How the model's variables are
          initialized can also affect it.
        \item \textbf{Asynchronous parallelism.} The order in which different
          parts of the model are updated can affect it.
        \item \textbf{Small evaluation sets.} If the evaluation set is too
          small, it may not be representative, producing uneven variations.
      \end{itemize}
      \vskip0.4em
      \footnotesize
      \gterm{Running an experiment multiple times helps confirm experimental
      results.}
    \end{tdmain}
    \begin{tdmargin}
      Your team should have a clear understanding of what exactly an experiment
      is, with a defined set of practices and artifacts: what the artifacts are,
      what coding practices apply, and what the standards are for
      reproducibility and tracking.
    \end{tdmargin}
  \end{withmargin}
\end{frame}

\begin{frame}{Wrong predictions, and stalled projects}
  \begin{withmargin}
    \begin{tdmain}
      \textbf{No real-world model is perfect.} Begin thinking early about how to
      deal with wrong predictions. A best-practices strategy encourages users to
      correctly label them: mail apps capture misclassified email by logging the
      mail users move into their spam folder, and the reverse.
      \vskip0.7em
      \textbf{Strategic fix for a stalled project.} You might need to
      \gterm{reframe} the problem. Maybe you initially wanted linear regression
      to predict a numeric value, but the data wasn't good enough. Further
      analysis might reveal the problem can be solved by predicting whether an
      example is above or below a specific value, reframing it as binary
      classification.
    \end{tdmain}
    \begin{tdmargin}
      \textbf{Technical fix.} Spend time diagnosing and analyzing wrong
      predictions. Isolate a few and diagnose the model's behavior in those
      instances. You might uncover problems with the architecture or the data.
      \par\medskip
      Maybe no solution is possible. Time-box your approach, stopping if you
      haven't made progress within the timeframe. However, if you have a strong
      problem statement, then it probably warrants a solution.
    \end{tdmargin}
  \end{withmargin}
  \keyterms{baseline}
\end{frame}

% ======================================================================
\sectionsubtitle{The goal is not one model. It is pipelines.}
\section[Pipelines]{ML Pipelines}

\begin{frame}{Why pipelines}
  \vfill
  \begin{takeaway}
    In production ML, the goal isn't to build a single model and deploy it.\\
    The goal is to build \alertbf{automated pipelines} for developing,\\
    testing, and deploying models over time.
  \end{takeaway}
  \vfill
  {\small As the world changes, trends in data shift, causing models in
  production to go stale. Without pipelines, replacing a stale model is an
  error-prone process: someone has to manually collect and process new data,
  train a new model, validate its quality, and finally deploy it.}
\end{frame}

\begin{frame}{The four pipelines}
  \begin{withmargin}
    \begin{tdmain}
      \textbf{Serving pipeline.} Delivers predictions. It exposes your model to
      the real world, receiving and processing the user's data, making a
      prediction, and delivering it.
      \vskip0.5em
      \textbf{Data pipeline.} Processes user data to create training and test
      datasets.
      \vskip0.5em
      \textbf{Training pipeline.} Trains models using the new training datasets
      from the data pipeline.
      \vskip0.5em
      \textbf{Validation pipeline.} Validates the trained model by comparing it
      with the production model using test datasets generated by the data
      pipeline.
    \end{tdmain}
    \begin{tdmargin}
      {\color{tdfg}\textbf{How they keep a fresh model in production.}}
      \par\medskip
      A model goes into production and the serving pipeline starts delivering
      predictions. The data pipeline immediately begins collecting data. Based
      on a schedule or trigger, the training and validation pipelines train and
      validate a new model. When validation confirms the new model isn't worse
      than the production model, the new model gets deployed. This process
      \gterm{repeats continuously}.
    \end{tdmargin}
  \end{withmargin}
\end{frame}

\begin{frame}{Model staleness and training frequency}
  \begin{withmargin}
    \begin{tdmain}
      Models tend to go stale almost immediately after they go into production.
      Their training datasets captured the state of the world a day ago, or in
      some cases, an hour ago. Inevitably the world has changed.
      \vskip0.7em
      Some models go stale faster than others. Models that recommend clothes
      typically go stale quickly because consumer preferences are notorious for
      changing frequently. Models that identify flowers might
      \gterm{never} go stale: a flower's identifying characteristics remain
      stable.
      \vskip0.5em
      \footnotesize
      A recommended best practice is to train and deploy new models on a
      \textbf{daily} basis. Just like regular software projects that have a
      daily build and release process.
    \end{tdmain}
    \begin{tdmargin}
      Establish a training frequency that reflects the nature of your data. If
      the data is dynamic, train often. If it's less dynamic, you might not need
      to train that often.
      \par\medskip
      Train models \textbf{before} they go stale. Early training provides a
      buffer to resolve potential issues, for example if the data or training
      pipeline fails, or the model quality is poor.
      \par\medskip
      Quick to go stale: spam, clothes recommendations, fraud detection. Not:
      bird species classification.
    \end{tdmargin}
  \end{withmargin}
\end{frame}

\begin{frame}{Online and offline predictions}
  \begin{withmargin}
    \begin{tdmain}
      \textbf{Online predictions} happen in real time, typically by sending a
      request to an online server and returning a prediction.
      \vskip0.6em
      \textbf{Offline predictions} are precomputed and cached. To serve a
      prediction, the app finds the cached prediction in the database and
      returns it. A subscription service might predict the churn rate for every
      subscriber and cache it; when the app needs it, for instance to incentive
      users about to churn, it just looks up the precomputed prediction.
      \vskip0.6em
      \footnotesize
      Predictions typically get \gterm{post-processed} before delivery: to
      remove toxic or biased content, boost more authoritative content, present
      a diversity of results, demote clickbait, or remove results for legal
      reasons.
    \end{tdmain}
    \begin{tdmargin}
      The infrastructure is different for online and offline predictions. This
      course refers to both as the serving pipeline.
      \par\medskip
      The feature engineering step is typically built \textbf{within} the model
      and not a separate, stand-alone process. The data processing code in the
      serving pipeline is often nearly identical to that in the data pipeline.
    \end{tdmargin}
  \end{withmargin}
\end{frame}

\begin{frame}{Logging predictions and capturing ground truth}
  \begin{withmargin}
    \begin{tdmain}
      The serving pipeline should incorporate a repository to log model
      predictions and, if possible, the ground truth. By aggregating
      predictions, you can monitor the general quality of your model and
      determine if it's starting to lose quality. Generally, the production
      model's predictions should have the \gterm{same average} as the labels
      from the training dataset.
      \vskip0.8em
      In some cases the ground truth only becomes available much later. If a
      weather app predicts the weather six weeks into the future, the ground
      truth won't be available for six weeks.
    \end{tdmain}
    \begin{tdmargin}
      A mail app implicitly captures user feedback when users move mail from
      their inbox to their spam folder. However, this only works when users
      correctly categorize their mail. When users leave spam in their inbox,
      because they know it's spam and never open it, that mail gets labeled
      "not spam" when it should be "spam".
      \par\medskip
      Always try to find ways to capture and record the ground truth, but be
      aware of the shortcomings in feedback mechanisms.
    \end{tdmargin}
  \end{withmargin}
\end{frame}

\begin{frame}{Versioned repositories}
  \begin{withmargin}
    \begin{tdmain}
      \footnotesize
      \begin{tabular}{@{}p{0.24\linewidth}p{0.70\linewidth}@{}}
        \toprule
        \textbf{Benefit} & \textbf{What it buys you}\\
        \midrule
        Reproducibility & Recreate and standardize model training environments
          and compare prediction quality among different models\\[0.3em]
        Compliance & Adhere to regulatory requirements for auditability and
          transparency\\[0.3em]
        Retention & Set data retention values for how long to store the data\\[0.3em]
        Access management & Manage who can access your data through fine-grained
          permissions\\[0.3em]
        Data integrity & Track and understand changes to datasets over time,
          making it easier to diagnose issues\\[0.3em]
        Discoverability & Make it easy for others to find your datasets,
          features, and models\\
        \bottomrule
      \end{tabular}
    \end{tdmain}
    \begin{tdmargin}
      Model repositories additionally give you tracking and evaluation of models
      in production, an easy \gterm{model release process} to review, approve,
      release, or roll back, and reproducibility for debugging by tracing a
      model's datasets and dependencies across deployments.
      \par\medskip
      Use \textbf{data cards} to structure information about datasets you share,
      and \textbf{model cards} to document a model's purpose, architecture,
      hardware requirements, and evaluation metrics.
    \end{tdmargin}
  \end{withmargin}
\end{frame}

\begin{frame}{Six challenges building pipelines}
  \begin{withmargin}
    \begin{tdmain}
      \footnotesize
      \begin{itemize}
        \item \textbf{Getting access to the data you need.} You might need to
          explain how the data will be used and clarify how personal
          identifiable information issues will be solved.
        \item \textbf{Getting the right features.} The features used in
          experimentation might not be available from the real-time data.
        \item \textbf{Understanding how the data is collected and represented.}
          Don't use data you're not confident in to train a model that might go
          to production.
        \item \textbf{Understanding the tradeoffs between effort, cost, and
          model quality.} A new feature might require a lot of effort and only
          slightly improve quality, or be easy to add but prohibitively
          expensive to store.
        \item \textbf{Getting compute.} TPU quota can be hard to get, and
          managing TPUs is complicated.
        \item \textbf{Finding the right golden dataset.} If the data changes
          frequently, getting golden datasets with consistent and accurate
          labels can be challenging.
      \end{itemize}
    \end{tdmain}
    \begin{tdmargin}
      \vfill
      Catching these problems during \gterm{experimentation} saves time. You
      don't want to develop the best features and model only to learn they're
      not viable in production.
      \vfill
    \end{tdmargin}
  \end{withmargin}
  \keyterms{feature engineering; ground truth; offline inference; online
    inference; prediction bias; test set}
\end{frame}

% ======================================================================
\sectionsubtitle{Compute, monitoring, and the deployment runbook.}
\section[Productionization]{Productionization}

\begin{frame}{Provisioning compute}
  \begin{withmargin}
    \begin{tdmain}
      Running ML pipelines requires compute resources: RAM, CPUs, and GPUs or
      TPUs. Depending on your use case, you might train and serve on different
      hardware. In general, it's common to \gterm{train on bigger hardware and
      serve on smaller}.
      \vskip0.7em
      \begin{itemize}
        \item Can you train on less expensive hardware?
        \item Would switching to different hardware boost performance?
        \item What size is the model and which hardware will optimize its
          performance?
        \item What hardware is ideal based on your model's architecture?
      \end{itemize}
    \end{tdmain}
    \begin{tdmargin}
      Switching hardware might make the model cheaper to run, but the
      engineering effort to migrate might outweigh the savings.
      \par\medskip
      You might allocate quota that pipelines \textbf{share}. In such cases,
      verify you have enough to run all your pipelines, and set up monitoring
      and alerting to prevent a single, errant pipeline from consuming all the
      quota.
      \par\medskip
      When estimating quota, factor in not only production pipelines but also
      ongoing experiments.
    \end{tdmargin}
  \end{withmargin}
\end{frame}

\begin{frame}{What to monitor, per pipeline}
  \begin{withmargin}
    \begin{tdmain}
      \footnotesize
      \begin{tabular}{@{}p{0.20\linewidth}p{0.74\linewidth}@{}}
        \toprule
        \textbf{Pipeline} & \textbf{Monitor}\\
        \midrule
        Serving & Skews or drifts in the serving data compared to the training
          data; skews or drifts in predictions; data type issues like missing or
          corrupted values; quota usage; model quality metrics\\[0.4em]
        Data & Skews and drifts in feature values; skews and drifts in label
          values; data type issues; quota usage rate; quota limit about to be
          reached\\[0.4em]
        Training & Training time; training failures; quota usage\\[0.4em]
        Validation & Skew or drift in the test datasets\\
        \bottomrule
      \end{tabular}
      \vskip0.4em
      Also monitor \gterm{latency}, how long it takes to deliver a prediction,
      and \gterm{outages}, whether the model has stopped delivering predictions
      at all.
    \end{tdmain}
    \begin{tdmargin}
      Good logging and monitoring practices help proactively identify issues and
      mitigate potential business impact. When issues do occur, alerts notify
      your team, and comprehensive logs facilitate diagnosing the root cause.
      \par\medskip
      Comparing the \textbf{changes} in proxy metrics is more insightful than
      their raw numbers.
    \end{tdmargin}
  \end{withmargin}
\end{frame}

\begin{frame}{Deploying a model}
  \begin{withmargin}
    \begin{tdmain}
      Document the following before you need them.
      \vskip0.5em
      \begin{itemize}
        \item Approvals required to begin deployment and increase the roll out.
        \item How to put a model into production.
        \item Where the model gets deployed, for example whether there are
          staging or canary environments.
        \item What to do if a deployment fails.
        \item How to rollback a model already in production.
      \end{itemize}
      \vskip0.5em
      \footnotesize
      Typically you deploy new models to a \gterm{subset of users} to check that
      the model is behaving as expected. If it is, continue. If it's not,
      rollback and begin diagnosing.
    \end{tdmain}
    \begin{tdmargin}
      After automating model training, automate validation and deployment.
      Automating deployments distributes responsibility and reduces the
      likelihood of a deployment being bottlenecked by a single person. It also
      reduces potential mistakes, increases efficiency and reliability, and
      enables on-call rotations and SRE support.
    \end{tdmargin}
  \end{withmargin}
\end{frame}

% ======================================================================
\sectionsubtitle{Fairness, privacy, transparency, safety.}
\section[Ethics]{AI and ML Ethics}

\begin{frame}{Your projects should benefit society}
  \begin{withmargin}
    \begin{tdmain}
      ML has the potential to transform society in many meaningful ways, either
      positively or negatively. It's critical to consider the ethical
      implications of your models and the systems they're a part of.
      \vskip0.8em
      Your ML projects \alertbf{should benefit society}. They shouldn't cause
      harm or be susceptible to misuse. They shouldn't perpetuate, reinforce, or
      exacerbate biases or prejudices. They shouldn't collect or use personal
      data irresponsibly.
    \end{tdmain}
    \begin{tdmargin}
      Beyond adhering to responsible AI principles, aim to develop systems that
      incorporate fairness, privacy, transparency, and safety.
    \end{tdmargin}
  \end{withmargin}
\end{frame}

\begin{frame}{Fairness}
  \begin{withmargin}
    \begin{tdmain}
      Models exhibit bias when their training data doesn't reflect the
      real-world population of their users, preserves biased decisions or
      outcomes, or uses features with more predictive power for certain groups.
      \vskip0.5em
      \footnotesize
      \begin{itemize}
        \item Identify underrepresented groups in evaluation datasets. You might
          need to oversample a subgroup to increase their presence.
        \item Use golden datasets to validate the model against fairness issues
          and detect implicit bias.
        \item Avoid including sensitive features, like gender or ethnicity.
        \item Avoid including features with little empirical or explanatory
          power, especially in sensitive contexts.
        \item Measure potential adverse impact on particular groups, and
          consider intentional bias correction techniques.
      \end{itemize}
    \end{tdmain}
    \begin{tdmargin}
      In a model for approving home loans, \gterm{don't include names} in the
      training data. Not only is an applicant's name irrelevant to the
      prediction task, but leaving such an irrelevant feature in the dataset
      also has the potential to create implicit bias or allocative harms. The
      model might correlate male names with a higher probability for repayment,
      or vice versa.
      \par\medskip
      The first step is verifying the training data accurately reflects the
      distribution of your users.
    \end{tdmargin}
  \end{withmargin}
\end{frame}

\begin{frame}{Privacy, transparency, safety}
  \begin{withmargin}
    \begin{tdmain}
      \textbf{Privacy.} Incorporate privacy design principles from the
      beginning. Be aware of and adhere to the European Union's Digital Markets
      Act for consent to share or use personal data, and EU GDPR laws. Remove
      all personally identifiable information from datasets and confirm your
      model and data repositories have the right permissions, for example
      \gterm{not world-readable}.
      \vskip0.6em
      \textbf{Transparency.} Be accountable to people. Make it easy for others
      to understand what your model does, how it does it, and why it does it.
      \vskip0.6em
      \textbf{Safety.} Design models to operate safely in adversarial
      conditions. Test your model with potential hostile inputs and check for
      potential failure conditions.
    \end{tdmain}
    \begin{tdmargin}
      Model cards provide a template to document your model and create
      transparency artifacts.
      \par\medskip
      Teams typically use specially designed datasets to test their models with
      inputs or conditions that caused the model to fail in the past.
      \par\medskip
      Always consider the broader social contexts your models operate within.
    \end{tdmargin}
  \end{withmargin}
  \slidesources{Mitchell2019}
\end{frame}

\begin{frame}{Tools change; the nature of the problem does not}
  \begin{withmargin}
    \begin{tdmain}
      ML development requires using a variety of constantly evolving tools and
      frameworks. New ML tools continue to emerge as ways to handle complex data
      types, advances in hardware, and techniques for orchestrating pipelines
      continue to develop.
      \vskip0.8em
      As a result, companies, organizations, and teams implement ML solutions
      using different tools and frameworks, which likely change over time. While
      common frameworks and best practices are emerging, keep in mind that the
      nature of your particular problem might require \gterm{custom solutions}
      in certain cases.
    \end{tdmain}
    \begin{tdmargin}
      That is why this track spends its time on the phases, the data, the
      metrics, and the pipelines rather than on any one framework.
    \end{tdmargin}
  \end{withmargin}
  \slidesources{GoogleManagingML2025}
\end{frame}
